\chapter{Some Facts about the Binomial and Poisson Distributions}

\section{Poisson Distribution}

\begin{fact}[Sum of Independent Poisson is Poisson]\label{Fact:poissonsum}
If $X_1\sim\mathrm{Poisson}(\lambda_1)$,
$X_2\sim\mathrm{Poisson}(\lambda_2)$,
and $X_1\perp X_2$, then 
$X_1+X_2\sim\Poisson(\lambda_1+\lambda_2)$.
\end{fact}

\begin{proof}
We compute the probability mass function of $X_1+X_2$:
\begin{align*}
   & \prob*{X_1+X_2=n}\\    
={}& \sum_{k=0}^n \prob*{X_1=k,X_2=n-k}\\
={}& \sum_{k=0}^n \prob*{X_1=k}\prob*{X_2=n-k}\\
={}& \sum_{k=0}^n e^{-\lambda_1}\frac{\lambda_1^k}{k!}
e^{-\lambda_2}\frac{\lambda_2^{n-k}}{(n-k)!}\\
={}& 
\frac{e^{-(\lambda_1+\lambda_2)}}{n!}
\sum_{k=0}^n \frac{n!}{k!(n-k)!}\lambda_1^k\lambda_2^{n-k}\\
={}&\frac{e^{-(\lambda_1+\lambda_2)}}{n!}
(\lambda_1+\lambda_2)^k\;.
\end{align*}
\end{proof}

\begin{fact}[Characteristic Function of Poisson Random Variables]
The characteristic function of $\mathrm{Poisso}(\lambda)$ distribution is the function $\phi:\real\to\complex$ given by 
\[
\phi(t) = \exp\left(\lambda(e^{it}-1)\right) \;.
\]
\end{fact}

\begin{proof}
Let $X$ denote a $\mathrm{Poisson}(\lambda)$ random variable. Then, for any $t\in\real$ we have
\begin{align*}
\phi(t) ={} & \E*{e^{itX}}\\
={} & \sum_{k=0}^\infty e^{itk} \, \prob*{X=k}\\
={} & \sum_{k=0}^\infty e^{itk} \, e^{-\lambda} \, \frac{\lambda^k}{k!}\\
={} & e^{-\lambda} \sum_{k=0}^\infty \frac{(\lambda e^{it})^k}{k!}\\
={} & e^{-\lambda} \, e^{\lambda e^{itk}}\\
={} & e^{\lambda (e^{it}-1)}\;.
\end{align*}
\end{proof}

\begin{fact}[Characteristic Function of Sum of Independent Random Variables]
If $X_1$ has ch.f. $\phi_1$, $X_2$ has ch.f. $\phi_2$, and $X_1\perp X_2$. Then
$X_1+X_2$ has ch.f. $\phi_1\times\phi_2$.
\end{fact}

\begin{fact}[Characteristic Function Characterize Distribution Function]
If $X$ and $Y$ have the same characteristic functions, then they must have the same distribution.    
\end{fact}

\begin{proof}[Alternative proof of Fact~\ref{Fact:poissonsum}]\;
\begin{itemize}
\item $X_1$ has characteristic function $\phi_1(t)=e^{\lambda_1(e^{it}-1)}$.
\item $X_2$ has characteristic function $\phi_2(t)=e^{\lambda_2(e^{it}-1)}$.
\item $X_1+X_2$ has characteristic function $\phi(t)=\phi_1(t)\phi_2(t)=e^{(\lambda_1+\lambda_2)(e^{it}-1)}$. 
\item Thus $X_1+X_2\sim\mathrm{Poisson}(\lambda_1+\lambda_2)$.
\end{itemize}
\end{proof}

\section{Binomial Distribution}

\begin{fact}[Sum of Independent Binomial]\label{Fact:binsum}
If $X\sim\Binomial(n,p)$, $Y\sim\Binomial(m,p)$, and $X\perp Y$, then $X+Y\sim\Binomial(n+m,p)$.
\end{fact}

\begin{proof}
We compute the probability mass function of $X+Y$:
\begin{align*}
& \prob*{X+Y=k}\\
={} & \sum_{i=0}^{k} \prob*{X=i,Y=k-i}\\
={} & \sum_{i=0}^{k} \prob*{X=i}\prob*{Y=k-i}\\
={} & \sum_{i=0}^{k} \binom{n}{i}p^i(1-p)^{n-i}
\binom{m}{k-i}p^{k-i}(1-p)^{m-k+i}\\
={} & p^k(1-p)^{n+m-k}\sum_{i=0}^{k} \binom{n}{i}\binom{m}{k-i}\\
={} & p^k(1-p)^{n+m-k}\binom{n+m}{k}\;.
\end{align*}
\end{proof}

\filbreak

\begin{fact}[Characteristic Function of Binomial]
The characteristic function of $\Binomial(n,p)$ distribution is the function $\phi:\real\to\complex$ given by  
Then    
\[
\phi(t) = \left(1-p+pe^{it}\right)^n\;.
\]
\end{fact}

\begin{proof}
\begin{align*}
 &  \E*{e^{itX}}\\
={}&\sum_{k=0}^n e^{itk} \prob*{X=k}\\
={}&\sum_{k=0}^n e^{itk} \binom{n}{k} p^k (1-p)^{n-k}\\
={}&\sum_{k=0}^n \binom{n}{k} (e^{it}p)^k (1-p)^{n-k}\\
={}&(1-p+pe^{it})^n\;.
\end{align*}
\end{proof}


\begin{proof}[Alternative proof of Fact~\ref{Fact:binsum}]
\;
\begin{itemize}
\item Characteristic function of $X$: $\phi_1(t)=(e^{it}p+1-p)^n$.
\item Characteristic function of $Y$: $\phi_2(t)=(e^{it}p+1-p)^m$.
\item Characteristic function of $X+Y$: $\phi(t)=\phi_1(t) \phi_2(t) = (e^{it}p+1-p)^{n+m}$.
\item So $X+Y\sim\Binomial(n+m,p)$.
\end{itemize}
\end{proof}

\filbreak

\section{A Connection Between Poisson and Binomial}

\begin{fact}[A Connection Between Poisson and Binomial]
If $X_1\sim\Poisson(\lambda_1)$, $X_2\sim\Poisson(\lambda_2)$, and $X_1\perp X_2$, then the conditional distribution of $X_1$ given $X_1+X_2=n$ is $\Binomial(n,\lambda_1/(\lambda_1+\lambda_2))$.
\end{fact}

\begin{proof}
We compute the conditional probability mass function of $X_1$ given $X_1+X_2$:
\begingroup
\addtolength{\jot}{1em}
\begin{align*}
    & \prob*{X_1=k\given X_1+X_2=n}\\
={} & \frac{\prob*{X_1=k,X_1+X_2=n}}{\prob*{X_1+X_2=n}}\\
={} & \frac{\prob*{X_1=k,X_2=n-k}}{\prob*{X_1+X_2=n}}\\
={} & \frac{\prob*{X_1=k}\prob*{X_2=n-k}}{\prob*{X_1+X_2=n}}\\
={} & \frac{\displaystyle
e^{-\lambda_1}\frac{\lambda_1^{k  }}{k!    } 
e^{-\lambda_2}\frac{\lambda_2^{n-k}}{(n-k)!}
}
{\displaystyle
e^{-(\lambda_1+\lambda_2)}\frac{(\lambda_1+\lambda_2)^n}{n!}
} \\
={} & \frac{n!}{k!(n-k)!}\frac{\lambda_1^k\lambda_2^{n-k}}{(\lambda_1+\lambda_2)^{n}} \\
={} & \binom{n}{k}\Bigl(\frac{\lambda_1}{\lambda_1+\lambda_2}\Bigr)^k
\Bigl(\frac{\lambda_2}{\lambda_1+\lambda_2}\Bigr)^{n-k}\\
={} & \binom{n}{k}\Bigl(\frac{\lambda_1}{\lambda_1+\lambda_2}\Bigr)^k
\Bigl(1-\frac{\lambda_1}{\lambda_1+\lambda_2}\Bigr)^{n-k}\;.
\end{align*}    
\endgroup
\end{proof}

\chapter{Combinatorial Identities}

\begin{fact}[A Combinatorial Identity]
For non-negative integers $n$, $m$, $k$
\[
\sum_{i=0}^{k} \binom{n}{i}\binom{m}{k-i} = \binom{n+m}{k} \;.
\]
\end{fact}

\begin{fact}[Another Combinatorial Identity]
\[
\sum_{m=0}^{N} \binom{2m}{m}\binom{2N-2m}{N-m} = 2^{2N}\;.
\]    
\end{fact}
